\documentclass[10pt,a4paper]{amsart}
\usepackage[margin=19mm]{geometry}
\usepackage{amsmath,amssymb,mathtools}
\usepackage[scr=rsfs,cal=euler]{mathalfa}
\usepackage{xfrac}
\usepackage[unicode,hidelinks,bookmarks=false]{hyperref}
\newcommand{\ie}{i.e.\ }
\newcommand{\cf}{cf.\ }
\newcommand{\resp}{respectively\ }
\newcommand{\Subsection}{\subsection}
\providecommand{\nobreakdash}{\nobreak\hbox{-}}
\setlength{\emergencystretch}{2em}
\multlinegap0pt
\usepackage{amssymb}
\usepackage{mathtools}
\usepackage[normalem]{ulem}
\usepackage[shortlabels]{enumitem}
\usepackage{xcolor}
\usepackage{tikz}
\newtheorem{proposition}{Proposition}
\newcommand{\B}{\mathcal{B}}
\newcommand{\Z}{\mathbb{Z}}
\newcommand{\cH}{\mathcal{H}}
\newcommand{\cK}{\mathcal{K}}
\newcommand{\fA}{\mathfrak{A}}
\newcommand{\fB}{\mathfrak{B}}
\newcommand{\fH}{\mathfrak{H}}
\newcommand{\fK}{\mathfrak{K}}
\newcommand{\ima}{\mathrm{im}}
\title{Computing Connection Matrices of Conley Complexes \\ via Algebraic Morse Theory}
\author{\'Alvaro Torras-Casas, Ka Man Yim, Ulrich Pennig}
\date{Source: arXiv:2503.09301v2 (CC BY 4.0)}
\begin{document}
\maketitle

\noindent\emph{Source and licence.} Computing Connection Matrices of Conley Complexes \\ via Algebraic Morse Theory. Version arXiv:2503.09301v2 (CC BY 4.0), distributed by arXiv under the Creative Commons Attribution 4.0 International licence, http://creativecommons.org/licenses/by/4.0/, as recorded in the arXiv licence metadata for this exact version and shown on the abstract page https://arxiv.org/abs/2503.09301v2. Full licence notice: \texttt{licenses/arxiv-2503.09301-CC-BY-4.0.txt}.\par\smallskip

\noindent\textit{Editorial scope.} Counted unit: the article's proposition prop:splittings-finite-dimensional (source0.tex lines 1378--1388). The article's complete original proof of it is delivered with it (source0.tex lines 1389--1422, one \texttt{proof} environment of 34 lines). No mathematics is written, repaired or re-derived: the statement and the proof are the article's own lines, in the article's own notation, and no retained line is substituted or re-flowed.  No further source-local context is needed: every cross-reference of every retained line resolves inside the two delivered spans.  Nothing was omitted from any retained span, so the package carries no omission record.  No retained line contains a citation, so the wrapper carries no bibliography and no bibliography entry.  No retained line contains a figure environment, an image-inclusion command or drawing code, so no image is delivered and the package declares no asset dependency.  Editorial formatting only, in addition: an AMS \texttt{amsart} preamble that restates the article's own declarations and macros used by the retained lines, namely the environments \texttt{proposition} on separate counters, together with the standard packages those lines need.  The retained environments print in this document as Proposition 1 in order of appearance, whereas the article prints the same items with the numbering of its own full document.  The complete native source of the article as distributed in the arXiv e-print for this exact version is bundled unchanged as \texttt{sources/174634/source0.tex}.
\begin{proposition}\label{prop:splittings-finite-dimensional}
    Let $(C,d)$ be a chain complex over a field $k$ which is bounded and whose groups are finite dimensional.
    Given a choice of bases $\fA_n$ for each $C_n$, there are split isomorphisms $\zeta_n\colon C_n\rightarrow H_n\oplus B_n\oplus K_n$ 
    for all $n\in \Z$, such that $\cH_n = \zeta_{n}^{-1}(H_n)$, $\B_n = \zeta_n^{-1}(B_n)$, and $\cK_n = \zeta_{n}^{-1}(K_n)$ and also 
\begin{itemize}
    \item $B_{n-1} = d_n(\cK_{n}) = \B_{n-1}$, and 
    \item $Z_n = \ker(d_n) =  \zeta^{-1}_n(H_n \oplus B_n) = \cH_n \oplus \B_n$
\end{itemize}
for all $n\in \Z$. 
In particular, notice that $d_n$ sends $\cK_n$ isomorphically onto $\B_{n-1}$.
\end{proposition}

\begin{proof}
The definition of $\zeta_n$ is given above for indices $0\leq n \leq N$; where $N\geq 0$ is the minimum positive integer such that $C_r=0$ for all $r>N$.
On the other hand, for indices $m \in \Z$ such that $m < 0$ or $m>N$, since $C_m=0$, the definition of $\zeta_m$ is trvial.
Now, by definition, it follows that $\cH_n=(\zeta^\cH_n)^{-1}(H_n)=\zeta_n^{-1}(H_n)$ as well as the other two equalites for inverse images. 

Next, for $0\leq n \leq N$, we show that  $\zeta_n=\zeta^\cH_n\oplus \zeta^\B_n\oplus \zeta^\cK_n$ is an isomorphism by proving that each component is also an isomorphism.
To start,  
\[
\B_n = \langle \fB_n\rangle = \langle c(R_{n+1})\rangle = \ima (D_n) = \ima(d_n) = B_n
\]
which implies that $\zeta^\B_n$ is an isomorphism. 
This also implies $B_n = \langle \fB_n\rangle$ and using that 
\[
\fH_n \sqcup \fB_n =\{ v \in c(T_n) \mid \mbox{ such that } D_nv=0\}.
\]
it follows that, vectors from $c(T_n)$ being linearly independent,
\[
\cH_n \oplus B_n = \langle\fH_n\rangle \oplus \langle \fB_n\rangle = 
\langle\fH_n \sqcup \fB_n\rangle = 
\ker(D_n)=\ker(d_n)=Z_n
\]
and so
$
\ima(\zeta^\cH_n)= (\cH_n \oplus B_n)/B_n = Z_n/B_n=H_n
$.
Thus, $\zeta^\cH_n$ is an isomorphism.
Finally, 
\[
\zeta^\cK_n(\cK_n)=d_n(\cK_n) = d_n(\langle\fK_n \rangle) = d_n(\langle c(T_n)\rangle) = d_n(C_n) = \ima(d_n)=B_{n-1}=K_n\ ,
\]
and since $\ker(d_n)\cap \cK_n=0$, it follows that $\zeta^\cK_n$ is an isomorphism.

The other properties follow from the argument above and the definition of $\zeta_n$.
\end{proof}

\end{document}
